% SOURCE_STATEMENT: "Does every finite-index subfactor admit a counterpart in conformal field theory, and can such a conformal net or vertex operator algebra be constructed directly from the subfactor data?"
% Preparation and final self-audit: CONTRIBUTING.md prompts C and D; source/audit artifacts retained in the local physics_statements_2026_10_11 batch.
% !TeX program = lualatex
\documentclass[11pt]{article}
\usepackage[margin=25mm]{geometry}
\usepackage{fontspec}
\setmainfont{Latin Modern Roman}
\usepackage{amsmath,amssymb,mathtools}
\usepackage{unicode-math}
\setmathfont{Latin Modern Math}
\usepackage{xurl,hyperref}
\hypersetup{colorlinks=true,urlcolor=blue}
\setcounter{secnumdepth}{0}
\setlength{\parindent}{0pt}
\setlength{\parskip}{5pt}
\setlength{\emergencystretch}{3em}
\newcommand{\sref}[2]{\hyperlink{source:#1:#2}{[S#2]}}
\newcommand{\cataloguescope}[1]{\paragraph{Recorded target.}#1}
\newcommand{\R}{\mathbb R}
\newcommand{\C}{\mathbb C}
\newcommand{\N}{\mathbb N}
\newcommand{\Z}{\mathbb Z}
\newcommand{\Q}{\mathbb Q}
\newcommand{\F}{\mathbb F}
\newcommand{\eps}{\varepsilon}
\begin{document}
\section{Conformal Field Theory Realizations of Subfactors}\label{statement}

\subsection{Definitions and mathematical statement}

\paragraph{Shared notation.}$\mathbb N=\{1,2,\ldots\}$; the standard number fields have their usual meanings. All additional conventions are stated locally.


\paragraph{Definitions and conventions.}
A subfactor is an inclusion $N\subset M$ of von Neumann algebras with trivial centers. Its index and standard invariant encode the size and higher relative-commutant structure of the inclusion. In an operator-algebraic conformal field theory, interval or local observable algebras and sectors can produce subfactor inclusions. To say a subfactor arises in conformal field theory requires identifying the inclusion or invariant that the construction realizes. Finite index restricts the inclusion to the finite Jones-index case. A conformal net is an interval-indexed family of local observable algebras with conformal covariance, while a vertex operator algebra is an algebraic chiral-field structure. A direct construction is to recover such a theory from the given subfactor data, not merely exhibit a categorical analogy.

\paragraph{Target retained from the canonical source catalogue.}
\paragraph{Catalogue wording retained for traceability.}
Does every finite-index subfactor admit a counterpart in conformal field theory, and can such a conformal net or vertex operator algebra be constructed directly from the subfactor data?

\paragraph{Source scope and qualification.}
This is related to ID 30002833 and contains the Asaeda--Haagerup instance, ID 30002840. The finite-index and constructive wording prevents asserting exact equivalence of these extracted records without matching the original realization notions.
\subsection{Short English statement}
Realize every finite-index subfactor in conformal theory and seek a direct construction from its data.
\subsection{Sources}
\begin{itemize}
\item[\hypertarget{source:statement:1}{S1}] ULAM.AI and the UnsolvedMath Contributors, UnsolvedMath, catalogue entry. Catalogue curation under CC BY 4.0; definitions and English exposition added locally.
\url{https://huggingface.co/datasets/ulamai/UnsolvedMath}.
\item[\hypertarget{source:statement:2}{S2}] Subfactors and Applications (2019), Oberwolfach Reports, DOI 10.4171/owr/2019/49, pp. 3053--3140. Primary report accessed 2026-10-11; relevant PDF pages 9.
\url{https://doi.org/10.4171/owr/2019/49}.
\item[\hypertarget{source:statement:3}{S3}] Jones, Does every subfactor arise in conformal field theory (catalogue-listed context; not independently consulted in this preparation).
\url{https://ems.press/content/serial-article-files/46563?nt=1}.
\item[\hypertarget{source:statement:4}{S4}] Bischoff, CFT realization of quantum doubles for index below four (catalogue-listed context; not independently consulted in this preparation).
\url{https://arxiv.org/abs/1506.02606}.
\end{itemize}
\label{end:statement}
\end{document}
